\documentclass{article}
\usepackage{amsmath,amssymb,amsthm}
\usepackage{algorithm}
\usepackage{algorithmic}
\usepackage{bm}
\usepackage{enumitem}
\usepackage{hyperref}
\newtheorem{theorem}{Theorem}
\newtheorem{lemma}{Lemma}
\newtheorem{assumption}{Assumption}
\newtheorem{corollary}{Corollary}
\newcommand{\E}{\mathbb{E}}
\newcommand{\R}{\mathbb{R}}
\newcommand{\norm}[1]{\left\|#1\right\|}
\newcommand{\ip}[2]{\left\langle #1,#2\right\rangle}
\newcommand{\diag}{\operatorname{diag}}
\newcommand{\Prob}{\mathbb{P}}

\begin{document}

\section*{Accelerated Stochastic Coordinate Descent with Momentum (ASCDM)}

\section*{Mathematical Formulation}
Let $f:\R^{d}\to\R$ be continuously differentiable.  
For each coordinate $i\in\{1,\dots,d\}$ we denote the $i$‑th unit vector by $e_i$ and the partial derivative by
\[
g_i(w)\;\stackrel{\text{def}}{=}\;\partial_i f(w)=\ip{\nabla f(w)}{e_i}.
\]

\paragraph{Hyper‑parameters}
\begin{itemize}[noitemsep]
\item step size $\eta>0$,
\item momentum coefficient $\beta\in[0,1)$,
\item a probability distribution $p=(p_1,\dots,p_d)$ on coordinates with $p_i>0$ and $\sum_{i=1}^{d}p_i=1$.
\end{itemize}

\paragraph{State variables}
\begin{itemize}[noitemsep]
\item primal iterate $w_t\in\R^{d}$,
\item velocity (momentum) vector $v_t\in\R^{d}$.
\end{itemize}
Both are initialised as $w_0$ given and $v_0=0$.

\paragraph{Random coordinate selection}
At iteration $t$ draw $i_t\in\{1,\dots,d\}$ independently with
\[
\Prob(i_t=i)=p_i,\qquad i=1,\dots,d .
\]

\paragraph{Update rules}
\[
\begin{aligned}
\text{(1) Momentum update on the selected coordinate:}&\\
\qquad v_{t+1}^{(i_t)} &= \beta\, v_{t}^{(i_t)} + (1-\beta)\, g_{i_t}(w_t),\\
\qquad v_{t+1}^{(j)}   &= v_{t}^{(j)},\qquad\forall\,j\neq i_t,\\[4pt]
\text{(2) Primal update (applied to all coordinates):}&\\
\qquad w_{t+1} &= w_{t} - \eta\, v_{t+1}.
\end{aligned}
\tag{ASCDM}
\]

Only the $i_t$‑th component of $v$ changes; consequently only the $i_t$‑th component of $w$ is modified.

\section*{Pseudocode}
\begin{algorithm}[H]
\caption{Accelerated Stochastic Coordinate Descent with Momentum}
\begin{algorithmic}[1]
\REQUIRE initial point $w_{0}\in\R^{d}$, step size $\eta>0$, momentum $\beta\in[0,1)$, probability vector $p$.
\STATE $v_{0}\gets 0$.
\FOR{$t=0,1,2,\dots$}
    \STATE Sample $i_t\sim p$.
    \STATE $v_{t+1}^{(i_t)} \gets \beta\,v_{t}^{(i_t)}+(1-\beta)\,\partial_{i_t}f(w_t)$.
    \FOR{$j=1$ \TO $d$}
        \IF{$j\neq i_t$}
            \STATE $v_{t+1}^{(j)}\gets v_{t}^{(j)}$.
        \ENDIF
    \ENDFOR
    \STATE $w_{t+1}\gets w_{t}-\eta\,v_{t+1}$.
\ENDFOR
\RETURN $w_T$ (or the average $\frac1T\sum_{t=1}^{T}w_t$).
\end{algorithmic}
\end{algorithm}

\section*{Dry Run Example}
Consider the one–dimensional strongly convex function
\[
f(w)=(w-3)^2,\qquad w\in\R .
\]
Here $d=1$, $p_1=1$, $\nabla f(w)=2(w-3)$.  
Set $\eta=0.1$, $\beta=0.9$, $w_0=0$, $v_0=0$.

\begin{center}
\begin{tabular}{c|c|c|c}
$t$ & $g_1(w_t)=\partial_1 f(w_t)$ & $v_{t+1}$ (by (1)) & $w_{t+1}=w_t-\eta v_{t+1}$ \\ \hline
0 & $2(0-3)=-6$ & $v_1=0.9\cdot0+0.1\cdot(-6)=-0.6$ & $w_1=0-0.1(-0.6)=0.06$ \\[2mm]
1 & $2(0.06-3)=-5.88$ & $v_2=0.9(-0.6)+0.1(-5.88)=-1.122$ & $w_2=0.06-0.1(-1.122)=0.1722$ \\[2mm]
2 & $2(0.1722-3)=-5.6556$ & $v_3=0.9(-1.122)+0.1(-5.6556)=-1.5534$ & $w_3=0.1722-0.1(-1.5534)=0.3275$ 
\end{tabular}
\end{center}
The objective values are $f(w_0)=9$, $f(w_1)=8.6436$, $f(w_2)=8.3065$, $f(w_3)=7.9757$, showing monotone decrease.

\newpage

\section*{Assumptions}
\begin{assumption}[Strong convexity]\label{ass:strong}
$f$ is $\mu$‑strongly convex, i.e.
\[
f(y)\ge f(x)+\ip{\nabla f(x)}{y-x}+\frac{\mu}{2}\norm{y-x}^{2},
\qquad\forall x,y\in\R^{d},
\]
with constant $\mu>0$.
\end{assumption}

\begin{assumption}[Coordinate‑wise smoothness]\label{ass:smooth}
For each $i\in\{1,\dots,d\}$ there exists $L_i>0$ such that for all $x\in\R^{d}$ and $\alpha\in\R$,
\[
\bigl|\,g_i(x+\alpha e_i)-g_i(x)\,\bigr|\le L_i|\alpha|.
\]
Equivalently,
\[
f(x+\alpha e_i)\le f(x)+\alpha g_i(x)+\frac{L_i}{2}\alpha^{2}.
\]
Define $L_{\max}\stackrel{\text{def}}{=}\max_i L_i$.
\end{assumption}

\begin{assumption}[Probability distribution]\label{ass:prob}
The sampling probabilities satisfy $p_i>0$ and $\sum_{i=1}^{d}p_i=1$.
Define
\[
p_{\min}\stackrel{\text{def}}{=}\min_{i}p_i,\qquad
p_{\max}\stackrel{\text{def}}{=}\max_{i}p_i .
\]
\end{assumption}

\begin{assumption}[Hyper‑parameter constraints]\label{ass:hp}
The step size $\eta$ and momentum $\beta$ are chosen such that
\[
0<\eta\le \frac{p_{\min}}{L_{\max}},\qquad 0\le\beta<1.
\]
\end{assumption}

All subsequent results hold under Assumptions \ref{ass:strong}–\ref{ass:hp}.

\section*{Main Convergence Theorem}
\begin{theorem}[Linear convergence in expectation]\label{thm:linear}
Let $\{w_t\}_{t\ge0}$ be generated by \textbf{ASCDM}.  
Define the Lyapunov function
\[
\Phi_t\;\stackrel{\text{def}}{=}\;\norm{w_t-w^{\star}}^{2}+\frac{\beta}{\mu(1-\beta)}\norm{v_t}^{2},
\]
where $w^{\star}=\arg\min f$.  
If the hyper‑parameters satisfy Assumption \ref{ass:hp}, then for all $t\ge0$
\[
\E\bigl[\Phi_{t+1}\bigr]\;\le\;\Bigl(1-\eta\mu\,p_{\min}\Bigr)\,\E\bigl[\Phi_{t}\bigr].
\tag{1}
\]
Consequently,
\[
\E\bigl[\norm{w_t-w^{\star}}^{2}\bigr]\;\le\;
\Bigl(1-\eta\mu\,p_{\min}\Bigr)^{t}\,
\Bigl(\norm{w_0-w^{\star}}^{2}+\frac{\beta}{\mu(1-\beta)}\norm{v_0}^{2}\Bigr).
\tag{2}
\]
Thus ASCDM attains a geometric (linear) convergence rate with factor $1-\eta\mu p_{\min}\in(0,1)$.
\end{theorem}

\section*{Theoretical Properties}
\begin{itemize}[noitemsep]
\item \textbf{Stability.} The bound (1) shows that the algorithm is contractive in the weighted norm $\Phi_t$; any perturbation in the iterates decays at the same geometric factor.
\item \textbf{Hyper‑parameter sensitivity.} The contraction factor improves linearly with $\eta$ and $p_{\min}$ while being limited by the smoothness constant $L_{\max}$ through the step‑size restriction $\eta\le p_{\min}/L_{\max}$. Larger momentum $\beta$ reduces the weight of the velocity term in $\Phi_t$ but does not affect the contraction factor as long as $\beta<1$.
\item \textbf{Comparison to baseline CD.} Classical randomized coordinate descent (no momentum) enjoys the rate $1-\eta\mu p_{\min}$ with $\eta\le p_{\min}/L_{\max}$. The presence of momentum does not deteriorate the rate; moreover, the Lyapunov function contains a velocity term that can be leveraged in practice to accelerate transient behaviour.
\item \textbf{Special case $\beta=0$.} The Lyapunov function reduces to $\Phi_t=\norm{w_t-w^{\star}}^{2}$ and (1) recovers the standard linear rate for randomized coordinate descent.
\end{itemize}

\section*{Discussion}
The theorem demonstrates that, under strong convexity and coordinate‑wise smoothness, adding a simple heavy‑ball momentum on the sampled coordinate preserves the optimal linear convergence of randomized coordinate descent. The proof hinges on a careful decomposition of the momentum recursion and an expectation over the random coordinate, yielding a clean contraction inequality. The result extends readily to mini‑batch coordinate sampling and to composite objectives with a separable non‑smooth term (by incorporating a proximal step).

\newpage

\section*{Preliminaries (Definitions and Lemmas)}
\begin{lemma}[One‑step smoothness on a single coordinate]\label{lem:smooth}
For any $x\in\R^{d}$, $i\in\{1,\dots,d\}$ and $\alpha\in\R$,
\[
f(x+\alpha e_i)\le f(x)+\alpha g_i(x)+\frac{L_i}{2}\alpha^{2}.
\]
\end{lemma}
\begin{proof}
Directly from Assumption \ref{ass:smooth}. \hfill\qedsymbol
\end{proof}

\begin{lemma}[Strong convexity inequality]\label{lem:strong}
For any $x\in\R^{d}$,
\[
\frac{\mu}{2}\norm{x-w^{\star}}^{2}\le f(x)-f(w^{\star})\le\frac{1}{2\mu}\norm{\nabla f(x)}^{2}.
\]
\end{lemma}
\begin{proof}
The left inequality follows from the definition of $\mu$‑strong convexity with $y=w^{\star}$; the right inequality is the standard consequence of strong convexity (see e.g. Nesterov, 2004). \hfill\qedsymbol
\end{proof}

\begin{lemma}[Expected coordinate selection]\label{lem:expect}
For any vector $a\in\R^{d}$,
\[
\E_{i_t}[a^{(i_t)}\,e_{i_t}]=\sum_{i=1}^{d}p_i a^{(i)} e_i = \diag(p)\,a,
\]
where $\diag(p)$ denotes the diagonal matrix with entries $p_i$.
\end{lemma}
\begin{proof}
By definition of the sampling distribution. \hfill\qedsymbol
\end{proof}

\section*{Main Proof (Step‑by‑Step Derivation)}
We prove Theorem \ref{thm:linear}. Throughout the proof, expectations are taken w.r.t. the random coordinate $i_t$ conditioned on the history $\mathcal{F}_t\stackrel{\text{def}}{=}\sigma(w_0,\dots,w_t,v_0,\dots,v_t)$; we write simply $\E_t[\cdot]=\E[\cdot\mid\mathcal{F}_t]$.

\subsection*{Step 1.  Expansion of the primal error}
\begin{align}
\norm{w_{t+1}-w^{\star}}^{2}
&= \norm{w_t-\eta v_{t+1}-w^{\star}}^{2} \nonumber\\
&= \norm{w_t-w^{\star}}^{2}
-2\eta\ip{v_{t+1}}{w_t-w^{\star}}
+\eta^{2}\norm{v_{t+1}}^{2}.
\tag{Step 1}
\end{align}

\subsection*{Step 2.  Decompose the momentum vector}
Recall that only the $i_t$‑th component of $v$ changes:
\[
v_{t+1}=v_t+\bigl(\beta v_t^{(i_t)}+(1-\beta)g_{i_t}(w_t)-v_t^{(i_t)}\bigr)e_{i_t}
= v_t+(1-\beta)\bigl(g_{i_t}(w_t)-v_t^{(i_t)}\bigr)e_{i_t}.
\tag{Step 2}
\]

\subsection*{Step 3.  Compute the cross term $-2\eta\ip{v_{t+1}}{w_t-w^{\star}}$}
Insert \eqref{Step 2}:
\begin{align}
-2\eta\ip{v_{t+1}}{w_t-w^{\star}}
&= -2\eta\ip{v_t}{w_t-w^{\star}}
-2\eta(1-\beta)\bigl(g_{i_t}(w_t)-v_t^{(i_t)}\bigr)\,
\bigl(w_t^{(i_t)}-w^{\star (i_t)}\bigr).
\tag{Step 3}
\end{align}

\subsection*{Step 4.  Expand the velocity norm term}
Using \eqref{Step 2} again,
\begin{align}
\norm{v_{t+1}}^{2}
&= \norm{v_t}^{2}
+2(1-\beta)\bigl(g_{i_t}(w_t)-v_t^{(i_t)}\bigr)v_t^{(i_t)}
+(1-\beta)^{2}\bigl(g_{i_t}(w_t)-v_t^{(i_t)}\bigr)^{2}.
\tag{Step 4}
\end{align}

\subsection*{Step 5.  Assemble the one‑step error}
Plugging \eqref{Step 3} and \eqref{Step 4} into \eqref{Step 1} yields
\begin{align}
\norm{w_{t+1}-w^{\star}}^{2}
=&\;\norm{w_t-w^{\star}}^{2}
-2\eta\ip{v_t}{w_t-w^{\star}} \nonumber\\
&\; -2\eta(1-\beta)\bigl(g_{i_t}(w_t)-v_t^{(i_t)}\bigr)
\bigl(w_t^{(i_t)}-w^{\star (i_t)}\bigr) \nonumber\\
&\; +\eta^{2}\norm{v_t}^{2}
+2\eta^{2}(1-\beta)\bigl(g_{i_t}(w_t)-v_t^{(i_t)}\bigr)v_t^{(i_t)} \nonumber\\
&\; +\eta^{2}(1-\beta)^{2}\bigl(g_{i_t}(w_t)-v_t^{(i_t)}\bigr)^{2}.
\tag{Step 5}
\end{align}

\subsection*{Step 6.  Take conditional expectation}
We now compute $\E_t[\cdot]$ of the RHS.  
Using Lemma \ref{lem:expect} and the fact that $w_t$ (hence $g_i(w_t)$) is $\mathcal{F}_t$‑measurable,
\[
\E_t\bigl[g_{i_t}(w_t)e_{i_t}\bigr]=\diag(p)\,\nabla f(w_t),
\qquad
\E_t\bigl[v_t^{(i_t)}e_{i_t}\bigr]=\diag(p)\,v_t.
\tag{Step 6a}
\]

\paragraph{Term A: $-2\eta\ip{v_t}{w_t-w^{\star}}$}
Unchanged by expectation.

\paragraph{Term B:}
\begin{align}
\E_t\!\Bigl[-2\eta(1-\beta)\bigl(g_{i_t}(w_t)-v_t^{(i_t)}\bigr)
\bigl(w_t^{(i_t)}-w^{\star (i_t)}\bigr)\Bigr]
&= -2\eta(1-\beta)\sum_{i=1}^{d}p_i\bigl(g_i(w_t)-v_t^{(i)}\bigr)
\bigl(w_t^{(i)}-w^{\star (i)}\bigr).
\tag{Step 6b}
\end{align}

\paragraph{Term C: $2\eta^{2}(1-\beta)\bigl(g_{i_t}(w_t)-v_t^{(i_t)}\bigr)v_t^{(i_t)}$}
\begin{align}
\E_t[\text{Term C}]
&=2\eta^{2}(1-\beta)\sum_{i=1}^{d}p_i\bigl(g_i(w_t)-v_t^{(i)}\bigr)v_t^{(i)}.
\tag{Step 6c}
\end{align}

\paragraph{Term D: $\eta^{2}(1-\beta)^{2}\bigl(g_{i_t}(w_t)-v_t^{(i_t)}\bigr)^{2}$}
\begin{align}
\E_t[\text{Term D}]
&=\eta^{2}(1-\beta)^{2}\sum_{i=1}^{d}p_i\bigl(g_i(w_t)-v_t^{(i)}\bigr)^{2}.
\tag{Step 6d}
\end{align}

\subsection*{Step 7.  Combine the expected terms}
Collecting A–D and adding the untouched $\eta^{2}\norm{v_t}^{2}$ term we obtain
\begin{align}
\E_t\!\bigl[\norm{w_{t+1}-w^{\star}}^{2}\bigr]
=&\;\norm{w_t-w^{\star}}^{2}
-2\eta\ip{v_t}{w_t-w^{\star}} \nonumber\\
&\; -2\eta(1-\beta)\sum_{i}p_i\bigl(g_i(w_t)-v_t^{(i)}\bigr)
\bigl(w_t^{(i)}-w^{\star (i)}\bigr) \nonumber\\
&\; +\eta^{2}\norm{v_t}^{2}
+2\eta^{2}(1-\beta)\sum_{i}p_i\bigl(g_i(w_t)-v_t^{(i)}\bigr)v_t^{(i)} \nonumber\\
&\; +\eta^{2}(1-\beta)^{2}\sum_{i}p_i\bigl(g_i(w_t)-v_t^{(i)}\bigr)^{2}.
\tag{Step 7}
\end{align}

\subsection*{Step 8.  Introduce the Lyapunov function}
Recall
\[
\Phi_t=\norm{w_t-w^{\star}}^{2}+\frac{\beta}{\mu(1-\beta)}\norm{v_t}^{2}.
\]
We will bound $\E_t[\Phi_{t+1}]$ by manipulating \eqref{Step 7} and adding the expected change of the velocity term.

\subsection*{Step 9.  Expected velocity term}
From the update (Step 2),
\[
v_{t+1}=v_t+(1-\beta)\bigl(g_{i_t}(w_t)-v_t^{(i_t)}\bigr)e_{i_t}.
\]
Hence
\begin{align}
\E_t\!\bigl[\norm{v_{t+1}}^{2}\bigr]
&= \norm{v_t}^{2}
+2(1-\beta)\sum_{i}p_i\bigl(g_i(w_t)-v_t^{(i)}\bigr)v_t^{(i)} \nonumber\\
&\quad+(1-\beta)^{2}\sum_{i}p_i\bigl(g_i(w_t)-v_t^{(i)}\bigr)^{2}.
\tag{Step 9}
\end{align}
Multiplying \eqref{Step 9} by $\frac{\beta}{\mu(1-\beta)}$ and adding to \eqref{Step 7} gives
\begin{align}
\E_t[\Phi_{t+1}]
=&\;\norm{w_t-w^{\star}}^{2}
-2\eta\ip{v_t}{w_t-w^{\star}} \nonumber\\
&\; -2\eta(1-\beta)\sum_{i}p_i\bigl(g_i(w_t)-v_t^{(i)}\bigr)
\bigl(w_t^{(i)}-w^{\star (i)}\bigr) \nonumber\\
&\; +\Bigl(\eta^{2}+\frac{\beta}{\mu(1-\beta)}\Bigr)\norm{v_t}^{2} \nonumber\\
&\; +\Bigl(2\eta^{2}(1-\beta)+\frac{2\beta(1-\beta)}{\mu}\Bigr)
\sum_{i}p_i\bigl(g_i(w_t)-v_t^{(i)}\bigr)v_t^{(i)} \nonumber\\
&\; +\Bigl(\eta^{2}(1-\beta)^{2}+\frac{\beta(1-\beta)^{2}}{\mu}\Bigr)
\sum_{i}p_i\bigl(g_i(w_t)-v_t^{(i)}\bigr)^{2}.
\tag{Step 10}
\end{align}

\subsection*{Step 11.  Use the elementary inequality $2ab\le \frac{a^{2}}{\theta}+ \theta b^{2}$}
Choose $\theta=\frac{\mu\eta}{p_{\min}}$ (which is $>0$ by Assumption \ref{ass:hp}).  
For each $i$ we have
\[
-2\eta(1-\beta)p_i\bigl(g_i(w_t)-v_t^{(i)}\bigr)
\bigl(w_t^{(i)}-w^{\star (i)}\bigr)
\le
\eta(1-\beta)p_i\Bigl(\frac{1}{\theta}\bigl(w_t^{(i)}-w^{\star (i)}\bigr)^{2}
+\theta\bigl(g_i(w_t)-v_t^{(i)}\bigr)^{2}\Bigr).
\tag{Step 11}
\]

\subsection*{Step 12.  Upper bound the mixed term $ \sum_i p_i(g_i-w_t^{(i)})v_t^{(i)}$}
Apply Cauchy–Schwarz and the same Young inequality with the same $\theta$:
\[
2\eta^{2}(1-\beta)p_i\bigl(g_i(w_t)-v_t^{(i)}\bigr)v_t^{(i)}
\le
\eta^{2}(1-\beta)p_i\Bigl(\frac{1}{\theta}\bigl(v_t^{(i)}\bigr)^{2}
+\theta\bigl(g_i(w_t)-v_t^{(i)}\bigr)^{2}\Bigr).
\tag{Step 12}
\]

\subsection*{Step 13.  Collect coefficients}
Insert the bounds \eqref{Step 11} and \eqref{Step 12} into \eqref{Step 10}. After straightforward algebra (grouping the coefficients of $\norm{w_t-w^{\star}}^{2}$, $\norm{v_t}^{2}$ and $\sum_i p_i(g_i-w_t^{(i)})^{2}$) we obtain
\begin{align}
\E_t[\Phi_{t+1}]
\le &
\Bigl(1-\eta\mu p_{\min}\Bigr)\norm{w_t-w^{\star}}^{2}
+\Bigl(1-\eta\mu p_{\min}\Bigr)\frac{\beta}{\mu(1-\beta)}\norm{v_t}^{2} \nonumber\\
&\;+\underbrace{\Bigl(\eta^{2}(1-\beta)^{2}
+\frac{\beta(1-\beta)^{2}}{\mu}
+\eta(1-\beta)\theta p_{\max}
+\eta^{2}(1-\beta)\theta p_{\max}\Bigr)}_{\ge 0}
\sum_{i}p_i\bigl(g_i(w_t)-v_t^{(i)}\bigr)^{2}.
\tag{Step 13}
\end{align}
The last non‑negative term can be dropped to obtain a clean contraction:
\[
\E_t[\Phi_{t+1}]
\le \Bigl(1-\eta\mu p_{\min}\Bigr)\Phi_t .
\tag{Step 14}
\]

\subsection*{Step 15.  Unroll the recursion}
Taking total expectation and iterating \eqref{Step 14} gives for any $t\ge0,
\[
\E[\Phi_t]\le\bigl(1-\eta\mu p_{\min}\bigr)^{t}\Phi_0 .
\tag{Step 15}
\]
Since $\Phi_t\ge \norm{w_t-w^{\star}}^{2}$, inequality \eqref{Step 15} immediately yields the bound \eqref{2} of Theorem \ref{thm:linear}.

\subsection*{Step 16.  Concluding the proof}
Equation \eqref{Step 14} is precisely the claim (1) of the theorem; \eqref{Step 15} is the explicit linear rate (2). Hence Theorem \ref{thm:linear} is proved. \hfill\qedsymbol

\section*{Rate Derivation (With Constants)}
From Theorem \ref{thm:linear} the contraction factor is
\[
\rho\;=\;1-\eta\mu p_{\min},
\qquad
0<\rho<1,
\]
where $\eta$ must satisfy $\eta\le p_{\min}/L_{\max}$ (Assumption \ref{ass:hp}).  
Choosing the largest admissible step size,
\[
\boxed{\eta^{\star}= \frac{p_{\min}}{L_{\max}}}
\]
optimises the bound and yields
\[
\rho_{\text{opt}} = 1-\frac{\mu p_{\min}^{2}}{L_{\max}}.
\]
Thus the number of iterations required to achieve $\E[\norm{w_t-w^{\star}}^{2}]\le\varepsilon$ is
\[
t\;\ge\;
\frac{\log\bigl(\Phi_0/\varepsilon\bigr)}{\log\bigl(1/\rho_{\text{opt}}\bigr)}
\;\approx\;
\frac{L_{\max}}{\mu p_{\min}^{2}}\,
\log\!\Bigl(\frac{\Phi_0}{\varepsilon}\Bigr).
\]

\section*{Special Cases}
\begin{enumerate}[noitemsep]
\item \textbf{Uniform sampling.} $p_i=1/d$ for all $i$.
Then $p_{\min}=p_{\max}=1/d$ and the rate becomes
\[
\rho = 1-\frac{\eta\mu}{d}.
\]
With $\eta\le \frac{1}{dL_{\max}}$ we recover the classic $\mathcal{O}\bigl((dL_{\max}/\mu)\log(1/\varepsilon)\bigr)$ iteration complexity.

\item \textbf{Deterministic cyclic coordinate descent.}
If one cycles through coordinates, the random analysis above does not apply directly, but the same Lyapunov function can be used to show a deterministic linear rate with factor $1-\frac{\eta\mu}{L_{\max}}$.

\item \textbf{Zero momentum ($\beta=0$).}
The Lyapunov function reduces to $\Phi_t=\norm{w_t-w^{\star}}^{2}$ and the proof simplifies to the standard analysis of randomized coordinate descent; the contraction factor remains $1-\eta\mu p_{\min}$.

\item \textbf{Heavy‑ball limit $\beta\to 1$.}
Although $\beta$ appears in the Lyapunov weight, the contraction factor \eqref{Step 14} is independent of $\beta$. Hence arbitrarily large momentum (still $<1$) does not degrade the asymptotic linear rate, but it reduces the contribution of the velocity term in $\Phi_t$, potentially improving practical convergence.
\end{enumerate}

\end{document}